\documentclass[11pt,a4paper]{article}
\usepackage[margin=1in]{geometry}
\usepackage{hyperref}
\usepackage{enumitem}

\title{Strategic Management \\ Week 10 \\ Slide-by-Slide Notes (ELI5) \\ Unit 5}
\author{(Generated from }sm-week-10-slides-1.pdf\text{)}
\date{}

\begin{document}
\maketitle

\section*{How to study this file (ELI5)}
\begin{itemize}
  \item For every slide, write the 2-3 ``Exam write'' lines in your notebook.
  \item Then memorize the ELI5 explanation (it is the reason behind the model).
  \item Use the slide keywords only as hints; always add the ``why'' in your own words.
\end{itemize}

\section*{Slide 1 of 18 (Contents)}
\textbf{Key idea (from slide):} Unit 5 competitive advantage and strategy.\\
\textbf{ELI5:} This unit teaches how firms win against rivals.\\
\textbf{Exam write:} ``Unit 5 is about how to compete using competitive advantage and clear positioning.''

\section*{Slide 2 of 18 (Learning Objectives)}
\textbf{Key idea (from slide):} Learn competitive advantage, analyze cost and differentiation, critique with the strategy clock, and connect to industry evolution.\\
\textbf{ELI5:} You learn 3 tools: cost, differentiation, and how the market changes.\\
\textbf{Exam write:} ``I will combine cost/diff with the industry life-cycle to choose the right strategy.''

\section*{Slide 3 of 18 (Competitive Advantage)}
\textbf{Key idea (from slide):}
\begin{itemize}
  \item Competitive advantage = firm key features vs rivals.
  \item Competitive parity = advantage not strong/ not sustainable.
  \item Customer value and margin relate to cost, price, and perceived value.
\end{itemize}
\textbf{ELI5:} If customers feel you are better \emph{and} you earn margin (gap between price and cost), you win.\\
\textbf{Exam write:} ``Competitive advantage is sustained, meaningful, and tied to customer value + margin.''

\section*{Slide 4 of 18 (Competitive Strategy)}
\textbf{Key idea (from slide):} Offensive-defensive actions to create a defendable position and achieve better performance. Includes focus: cost leader, product differentiation, and stuck-in-the-middle.\\
\textbf{ELI5:} You pick one main ``win way'' and protect it so rivals cannot copy easily.\\
\textbf{Exam write:} ``A competitive strategy is the defendable position you build: cost leader or product differentiation (not the middle).''

\section*{Slide 5 of 18 (Cost Leadership Advantage)}
\textbf{Key idea (from slide):} Cost leadership advantage when the firm costs are lower than competitors/industry average for similar quality/features.\\
\textbf{ELI5:} You win because your product costs less to make (and you still sell at a normal price).\\
\textbf{Exam write:} ``Cost leadership means lower cost for similar quality, so margin is higher.''

\section*{Slide 6 of 18 (Drivers of Cost Leadership)}
\textbf{Key idea (from slide):} Economies of scale; learning economies; process/production innovation (re-engineering); product design standardization (design for manufacture); input cost advantages; capacity utilization; ``residual efficiency'' (x-inefficiency, motivation).\\
\textbf{ELI5:} Cost gets lower when you run better systems: big volume, faster learning, better processes, cheaper inputs, and less waste.\\
\textbf{Exam write:} ``Cost advantage comes from many cost levers: scale, learning, process, design, inputs, capacity, and fewer inefficiencies.''

\section*{Slide 7 of 18 (Adequacy of Cost Leadership)}
\textbf{Key idea (from slide):} Cost leadership fits when customers are price-sensitive; competition pushes price down; product is standardized/commodity; buyers have high bargaining power; economies of scale and learning matter.\\
\textbf{ELI5:} If everyone compares mostly by price, being cheaper is strong.\\
\textbf{Exam write:} ``Cost leadership is adequate when price competition and buyer power are high and products are standardized.''

\section*{Slide 8 of 18 (Risks of Cost Leadership)}
\textbf{Key idea (from slide):} Risks include constant monitoring costs, innovation can cancel the learning/experience effect, cost inflation, demand changes, excessive quality reduction, and imitation/segmented competitors.\\
\textbf{ELI5:} Cheap strategy can break if you keep spending to manage costs, if competitors innovate, or if you cut quality too much.\\
\textbf{Exam write:} ``Risks: monitoring expense, innovation switching the advantage, inflation, demand shifts, quality damage, and competitor imitation.''

\section*{Slide 9 of 18 (Product Differentiation Advantage)}
\textbf{Key idea (from slide):} Product differentiation when firm offers unique attributes perceived by customers as ``unique'' and that creates a price premium.\\
\textbf{ELI5:} You charge higher because customers believe you are better.\\
\textbf{Exam write:} ``Differentiation is real only when uniqueness creates price premium (higher willingness to pay).''

\section*{Slide 10 of 18 (Drivers of Product Differentiation)}
\textbf{Key idea (from slide):}
\begin{itemize}
  \item Product features and performance + complementary services + perceived quality + location.
  \item Market features: variety of customer needs, perception, intangibles.
  \item Firm features: ethics/CSR, signaling and reputation, legitimacy, ``Other Time (ASAP)''.
\end{itemize}
\textbf{ELI5:} You differentiate by making customers feel: ``this is special'' in features, services, trust, or speed.\\
\textbf{Exam write:} ``Differentiate with features + service + perception/reputation (and sometimes time/speed), so customers accept premium price.''

\section*{Slide 11 of 18 (Adequacy of Product Differentiation)}
\textbf{Key idea (from slide):} Works when buyers care about quality; few competitors use the same differentiation criterion; product features are hard to imitate; in tech-based industries, innovation is the success factor.\\
\textbf{ELI5:} Differentiation works when ``different'' is valuable and rivals cannot copy fast.\\
\textbf{Exam write:} ``Differentiation is adequate when quality matters, rivals are not matching your exact basis, and imitation barriers are high.''

\section*{Slide 12 of 18 (Risks of Product Differentiation)}
\textbf{Key idea (from slide):}
\begin{itemize}
  \item Too high price premium.
  \item Demand changes during economic crises (customers may buy less).
  \item Imitation by competitors (turning your product into ``luxury'').
  \item Competitors focused on single segments.
\end{itemize}
\textbf{ELI5:} If you charge too much, or if customers get poorer, premium collapses. Also, if rivals copy you, the ``unique'' feeling disappears.\\
\textbf{Exam write:} ``Risks: premium too expensive, demand drops in crises, rivals imitate, and segment-focused competitors out-position you.''

\section*{Slide 13 of 18 (Extending Porter’s Strategies: Core Critical Learning)}
\textbf{Key idea (from slide):} Key teaching points:
\begin{itemize}
  \item Costs not always lead to competitive advantage (Cost leader is not automatically best).
  \item Stuck in the middle: there is only one cost leader horizontally.
  \item Product differentiation is not automatically price premium (they are different things).
  \item Brand products can act like luxury.
\end{itemize}
\textbf{ELI5:} Not every low-cost idea works. If you are ``in the middle'', nobody trusts you. And differentiation is only powerful if it truly changes what customers are willing to pay.\\
\textbf{Exam write:} ``Do not confuse cost with advantage, and do not confuse differentiation with price premium automatically; check the mechanism.''

\section*{Slide 14 of 18 (The Strategy Clock)}
\textbf{Key idea (from slide):} ``Strategy clock'' introduced as a way to critique simple cost/diff positioning and explore what is ``new''.\\
\textbf{ELI5:} Instead of only ``cheap'' vs ``unique'', strategy clock uses price and perceived value.\\
\textbf{Exam write:} ``Use strategy clock to pick a more accurate price-value position than just cost vs differentiation.''

\section*{Slide 15 of 18 (Strategic Clock Illustration)}
\textbf{Key idea (from slide):} Identify eight competitive strategy positions using the air passenger transport industry case, including drivers of positions.\\
\textbf{ELI5:} Airlines compete with different combos of low price vs high perceived value (speed, comfort, service).\\
\textbf{Exam write:} ``In the exam, map the case to a clock position using both price level and perceived value evidence.''

\section*{Slide 16 of 18 (Industry Life-Cycle Model)}
\textbf{Key idea (from slide):} Each industry stage needs a different competitive strategy: the ``way of competing'' plus costs, innovation, and profitability (linked to Five Forces).\\
\textbf{ELI5:} Markets are not static. In early stage, you invest; in mature stage, you defend and improve efficiency.\\
\textbf{Exam write:} ``Choose competitive strategy based on the industry life-cycle stage, not only on current competitors.''

\section*{Slide 17 of 18 (Life-Cycle Details: Emergence, Growth, Maturity, Decline)}
\textbf{Key idea (from slide):}
\begin{itemize}
  \item Emergence: shape industry structure; timing to entry (leader/follower); manage risk and finances.
  \item Growth: initiatives, investments; improve positioning (cost/differentiation); detect turning point (shakeout) and stop investment when needed.
  \item Maturity: consolidate advantage and competitive position; re-orient firm scope to emerging industry.
  \item Decline: check if temporal or structural; choices: harvest; segmentation; rapid withdrawal.
\end{itemize}
\textbf{ELI5:} In every life stage, the ``best move'' changes. You invest early, defend later, and sometimes harvest or exit when growth stops.\\
\textbf{Exam write:} ``In your answer, name the life-cycle stage and state the matching strategy logic (invest/defend/harvest/exit).''

\section*{Slide 18 of 18 (Business Case 3: Wal-Mart in China)}
\textbf{Key idea (from slide):} Case questions:
\begin{itemize}
  \item Why Wal-Mart success in the US? Competitive advantages and sources.
  \item Strategies Wal-Mart China should consider going forward.
\end{itemize}
\textbf{ELI5:} Compare US advantages to China context, then pick cost/diff choices that fit China competition and customer needs.\\
\textbf{Exam write:} ``Answer: (1) identify advantages with evidence; (2) choose next strategies that match China market stage and value drivers.''

\end{document}

